\documentclass[10pt,a4paper]{amsart}
\usepackage[margin=19mm]{geometry}
\usepackage{amsmath,amssymb,mathtools}
\usepackage[scr=rsfs,cal=euler]{mathalfa}
\usepackage{xfrac}
\usepackage[unicode,hidelinks,bookmarks=false]{hyperref}
\newcommand{\ie}{i.e.\ }
\newcommand{\cf}{cf.\ }
\newcommand{\resp}{respectively\ }
\newcommand{\Subsection}{\subsection}
\providecommand{\nobreakdash}{\nobreak\hbox{-}}
\setlength{\emergencystretch}{2em}
\multlinegap0pt
\usepackage{tikz}
\usepackage{amssymb}
\usepackage{mathtools}
\usepackage{xcolor}
\usepackage[shortlabels]{enumitem}
\usepackage{tcolorbox}
\usepackage{float}
\usepackage{csquotes}
\newtheorem{proposition}{Proposition}
\newtheorem{lemma}{Lemma}
\newcommand{\parenum}{\setenumerate[1]{label=\textbf{(\arabic*)}}}
\title{A Gentle Introduction to CAT(0) Spaces}
\author{S{\o}ren Poulsen}
\date{Source: arXiv:2406.09883v2 (CC BY 4.0)}
\begin{document}
\maketitle

\noindent\emph{Source and licence.} A Gentle Introduction to CAT(0) Spaces. Version arXiv:2406.09883v2 (CC BY 4.0), distributed by arXiv under the Creative Commons Attribution 4.0 International licence, http://creativecommons.org/licenses/by/4.0/, as recorded in the arXiv licence metadata for this exact version and shown on the abstract page https://arxiv.org/abs/2406.09883v2. Full licence notice: \texttt{licenses/arxiv-2406.09883-CC-BY-4.0.txt}.\par\smallskip

\noindent\textit{Editorial scope.} Counted unit: the article's lemma lem:induced\_convex (source0.tex lines 1341--1349). The article's complete original proof of it is delivered with it (source0.tex lines 1350--1450, one \texttt{proof} environment of 101 lines). No mathematics is written, repaired or re-derived: the statement and the proof are the article's own lines, in the article's own notation, and no retained line is substituted or re-flowed.  Retained uncounted as the source-local context the proof needs: lines 711--718.  Nothing was omitted from any retained span, so the package carries no omission record.  No retained line contains a citation, so the wrapper carries no bibliography and no bibliography entry.  No retained line contains a figure environment, an image-inclusion command or drawing code, so no image is delivered and the package declares no asset dependency.  Editorial formatting only, in addition: an AMS \texttt{amsart} preamble that restates the article's own declarations and macros used by the retained lines, namely the environments \texttt{proposition}, \texttt{lemma} on separate counters, together with the standard packages those lines need.  The retained environments print in this document as Proposition 1, Lemma 1 in order of appearance, whereas the article prints the same items with the numbering of its own full document.  The complete native source of the article as distributed in the arXiv e-print for this exact version is bundled unchanged as \texttt{sources/160803/source0.tex}.
\begin{proposition}\label{prop:conv_mid}
Let \((X,d)\) be a metric space such that any \(x\in X\) is contained in a geodesically convex neighborhood \(U_x\subseteq X\). The following are equivalent:
\parenum
\begin{enumerate}
    \item Any geodesic triangle \(\Delta(a,b,c)\) in \(U_x\) satisfies \(d(m,m')\leq\frac{1}{2}d(b,c)\), where \(m\) and \(m'\) are the midpoints of \([a,b]\) and \([a,c]\) respectively.
    \item Each pair of geodesics \(\gamma_1:[0,1]\to U_x,\gamma_2:[0,1]\to U_x\) parametrized proportional to the arc length satisfy \(d(\gamma_1(t),\gamma_2(t))\leq (1-t)d(\gamma_1(0),\gamma_2(0))+td(\gamma_1(1),\gamma_2(1))\) for all \(t\in[0,1]\).
\end{enumerate}
\end{proposition}

\begin{lemma}\label{lem:induced_convex}
Let \((X,d)\) be a metric space with locally complete and locally convex metric \(d\), i.e. where each point \(x\in X\) is contained in a geodesically convex and complete neighborhood \(U_x\) in which the induced metric is convex (see Proposition \ref{prop:conv_mid}).
Let \(\sigma:[0,1]\to X\) be a local geodesic between two points \(x,y\in X\), and let \(\varepsilon>0\) be such that for every \(t\in[0,1]\), the induced metric on the closed ball \(\overline{B}(\sigma(t),2\varepsilon)\) is complete and convex.
Then:
\parenum\begin{enumerate}
    \item For every points \(\overline{x},\overline{y}\in X\) with the properties \(d(x,\overline{x})<\varepsilon\) and \(d(y,\overline{y})<\varepsilon\), there is a unique local geodesic \(\overline{\sigma}:[0,1]\to X\) connecting \(\overline{x}\) to \(\overline{y}\) such that the map \(t\mapsto d(\sigma(t),\overline{\sigma}(t))\) is convex;
    \item \[L(\overline{\sigma})\leq L(\sigma)+d(x,\overline{x})+d(y,\overline{y}).\]
\end{enumerate}
\end{lemma}

\begin{proof}
We start by assuming that such a \(\overline{\sigma}\) exists.
First we show uniqueness.
Let \(\sigma_1,\sigma_2:[0,1]\to X\) be local geodesics such that condition (1) holds for both $\sigma_1,\sigma_2$. Then we have that for any such two local geodesics \(\sigma_1,\sigma_2\) satisfying the conditions of part (1), the map \(t\mapsto d(\sigma_1(t),\sigma_2(t))\) is convex. Indeed, first notice that 
$$d(\sigma_1(t),\sigma(t))\leq (1-t)d(\sigma_1(0),\sigma(0))+td(\sigma_1(1),\sigma(1))\leq (1-t)\varepsilon+\varepsilon=\varepsilon,$$ 
for all \(t\in [0,1]\) (and similarly for \(\sigma_2\)).
For all \(t\in [0,1]\) we have \(\sigma_1(t),\sigma_2(t)\in B(\sigma(t),2\varepsilon)\), which is a convex neighborhood of \(\sigma(t)\).
For every \(t\in [0,1]\) there is then a \(\delta_t>0\) such that \(v\mapsto d(\sigma_1(v),\sigma_2(v))\) is convex on \([t-\delta_t,t+\delta_t]\).
In particular, \(t\mapsto d(\sigma_1(t),\sigma_2(t))\) is locally convex, and therefore convex by a general theorem.
Now, by convexity and since \(\sigma_1(0)=\sigma_2(0)=\overline{x}\) and \(\sigma_1(1)=\sigma_2(1)=\overline{y}\) we must then have \(0\leq d(\sigma_1(t),\sigma_2(t))\leq (1-t)d(\sigma_1(0),\sigma_2(0))+td(\sigma_1(1),\sigma_2(1))=0\), so \(\sigma_1=\sigma_2\).
This shows that if a local geodesic satisfying the conditions in part (1) exists, it must necessarily be unique.

\medskip
Next we show that this local geodesic $\overline{\sigma}$ from (1) will satisfy the inequality in part (2).

 By the first part of the proof it is enough to consider local geodesics \(\sigma_1,\sigma_2:[0,1]\to X\) with \(d(\sigma(t),\sigma_1(t))<\varepsilon\) and \(d(\sigma(t),\sigma_2(t))<\varepsilon\) for all \(t\in [0,1]\), but this time we also assume \(\sigma_1(0)=\sigma_2(0)\). Notice 
\(\sigma_1(t)\) and \(\sigma_2(t)\) are both in the convex neighborhood \(B(\sigma(t),2\varepsilon)\) when \(t\in [0,1]\). So again by the convexity proven above \(d(\sigma_1(t),\sigma_2(t))\leq (1-t)d(\sigma_1(0),\sigma_2(0))+td(\sigma_1(1),\sigma_2(1))=td(\sigma_1(1),\sigma_2(1))\).
As \(\sigma_1\) and \(\sigma_2\) are local geodesics, \(\sigma_1|_{[0,t]}\) and \(\sigma_2|_{[0,t]}\) are geodesics for sufficiently small \(t\in (0,1]\).
In particular, \(L(\sigma_1|_{[0,t]})=tL(\sigma_1)\) and \(L(\sigma_2|_{[0,t]})=d(\sigma_2(0),\sigma_2(t))\).
We now have
\begin{align*}
tL(\sigma_2)&=d(\sigma_2(0),\sigma_2(t))=d(\sigma_1(0),\sigma_2(t))\\
            &\leq d(\sigma_1(0),\sigma_1(t))+d(\sigma_1(t),\sigma_2(t))\\
            &\leq tL(\sigma_1)+td(\sigma_1(1),\sigma_2(1)),
\end{align*}
for sufficiently small \(t\in (0,1]\).
Dividing through by \(t\) gives us
\begin{equation}\label{eq:length_ineq}L(\sigma_2)\leq L(\sigma_1)+d(\sigma_1(1),\sigma_2(1)),\end{equation}
for local geodesics \(\sigma_1,\sigma_2\) with \(\sigma_1(0)=\sigma_2(0)\).

Now let \(x,y,\overline{x},\overline{y}\) and \(\overline{\sigma}\) be as in part (1) of this Lemma. Let also \(\overline{\sigma_1}:[0,1]\to X\) be a local geodesic satisfying the condition (1) of this Lemma and with the property that $\overline{\sigma_1}(0)=\overline{\sigma}(0)=\overline{x}$ and $\overline{\sigma_1}(1)=\sigma(1)=y$.

Above we showed that those local geodesics $\overline{\sigma}, \overline{\sigma_1}$ are unique if they exist. Existence will be shown later. Then we apply inequality \eqref{eq:length_ineq} to get 
\begin{equation}\label{eq:length_ineq1}
L(\overline{\sigma})\leq L(\overline{\sigma_1})+d(\overline{\sigma}(1)),\overline{\sigma_1}(1)=L(\overline{\sigma_1})+d(y,\overline{y}).
\end{equation}
We also have \(-\overline{\sigma_1}(0)=-\sigma(0)=y\), so we can again apply inequality \eqref{eq:length_ineq}, which gives us \begin{align}
L(-\overline{\sigma_1})&\leq L(-\sigma)+d(-\overline{\sigma_1}(1),-\sigma(1)) \notag\\
L(\overline{\sigma_1})&\leq L(\sigma)+d(\overline{x},x)\label{eq:length_ineq2}.
\end{align}
Inserting inequality \eqref{eq:length_ineq1} into inequality \eqref{eq:length_ineq2} finally gives us
\[L(\overline{\sigma})\leq L(\sigma)+d(x,\overline{x})+d(y,\overline{y})\]
as desired.

\medskip
All of the above assumes the existence of such a \(\overline{\sigma}\).
We will now show \(\overline{\sigma}\) always exists.
Let \(A>0\), and consider the following statement:
\begin{enumerate}
    \item[\(P(A)\)]
For all \(a,b\in [0,1]\) with \(0<b-a\leq A\) and for all \(\overline{p},\overline{q}\in X\) with \(d(\sigma(a),\overline{p})<\varepsilon\) and \(d(\sigma(b),\overline{q})<\varepsilon\) there exists a local geodesic \(\overline{\sigma}:[a,b]\to X\) such that \(\overline{\sigma}(a)=\overline{p},\overline{\sigma}(b)=\overline{q}\) and such that \(d(\sigma(t),\overline{\sigma}(t))<\varepsilon\) for all \(t\in [a,b]\).
\end{enumerate}


First we notice that $P(A)$  is true for \(A:=\varepsilon/L(\sigma)\). Indeed, we have 
\(d(\sigma(a),\sigma(b))=L(\sigma)(b-a)=L(\sigma|_{[a,b]})\leq\varepsilon\).
For all \(\overline{q}\in X\) with \(d(\sigma(b),\overline{q})<\varepsilon\) we then have by the triangle inequality \(d(\sigma(a),\overline{q})\leq d(\sigma(a),\sigma(b))+d(\sigma(b),\overline{q})<2\varepsilon\).
This shows that all \(\overline{p},\overline{q}\in X\) as in the claim are in the ball \(\overline{B}(\sigma(a),2\varepsilon)\), which by hypothesis is geodesically convex.
Therefore, there exists a local geodesic from \(\overline{p}\) to \(\overline{q}\) which satisfies \(d(\sigma(t),\overline{\sigma}(t))<\varepsilon\) by convexity of \(d\) on \(B(\sigma(a),2\varepsilon)\), proving \(P(\varepsilon/L(\sigma))\).

\medskip
Now we know that $P(A)$ is true for at least one positive value of \(A\). If we can now show that \(P(A)\) implies \(P(3A/2)\), then the set of \(A\)'s for which the statement is true is unbounded, i.e. it is true for all \(A>0\). Indeed, assume $P(A)$ is true for some \(A>0\).
Let \(a,b\in [0,1]\) be such that \(0<b-a<3A/2\).
Next, split the interval \([a,b]\) into three smaller intervals with endpoints \(a<a_1<b_1<b\) of the same length.
Notice that \(b_1-a<A\) and \(b-a_1<A\).
Let \(\overline{p},\overline{q}\in X\) be as in the statement \(P(3A/2)\).
Define \(p_0=\sigma(a_1),q_0=\sigma(b_1)\).
Now we assumed that \(P(A)\) is true, so once we define points \(p_{n-1}\in X\) and \(q_{n-1}\in X\) we know there are local geodesics \(\sigma_n:[a,b_1]\to X\) and local geodesics \(\sigma'_n:[a_1,b]\to X\) joining \(\overline{p}\) to \(q_{n-1}\) and joining \(p_{n-1}\) to \(\overline{q}\), respectively, such that \(d(\sigma(t),\sigma_n(t))<\varepsilon\) for \(t\in [a,b_1]\) and \(d(\sigma(t),\sigma'_n(t))<\varepsilon\) for \(t\in [a_1,b]\).
We then define \(p_n=\sigma_n(a_1)\) and \(q_n=\sigma'_n(b_1)\).
Because we chose \(\varepsilon>0\) to be small enough for the metric to be convex we have 
\begin{align*}
d(p_0,p_1)&=d(\sigma(a_1),\sigma_1(a_1))\\
&\leq \frac{1}{2}d(\sigma(a),\sigma_1(a))+\frac{1}{2}d(\sigma(b_1),\sigma_1(b_1))\\
&<\frac{1}{2}\varepsilon
\end{align*}
and
\begin{align*}
d(p_n,p_{n+1})&=d(\sigma_n(a_1),\sigma_{n+1}(a_1))\\
&\leq\frac{1}{2}d(\sigma_n(a),\sigma_{n+1}(a))+\frac{1}{2}d(\sigma_n(b_1),\sigma_{n+1}(b_1))\\
&=\frac{1}{2}d(q_n,q_{n-1}).
\end{align*}
Similarly we get \(d(q_0,q_1)<\varepsilon/2\) and \(d(q_n,q_{n+1})\leq \frac{1}{2}d(p_n,p_{n-1})\).
Recursive substitution then yields
\begin{align*}
d(p_n,p_{n+1})&<\frac{\varepsilon}{2^{n+1}}\\
d(q_n,q_{n+1})&<\frac{\varepsilon}{2^{n+1}},
\end{align*}
showing that \((p_n)_{n\geq 0}\) and \((q_n)_{n\geq 0}\) are Cauchy sequences and therefore converge to unique points \(p_\infty\) and \(q_\infty\), respectively.

As \(t\mapsto d(\sigma_n(t),\sigma_{n+1}(t))\) is convex, it is bounded from above by \(d(\sigma_n(b_1),\sigma_{n+1}(b_1))<\frac{\varepsilon}{2^{n+1}}\).
This means that \((\sigma_n(t))_{n\geq 0}\) also is a Cauchy sequence in the complete ball \(\overline{B}(\sigma(t),\varepsilon)\) for all \(t\in [a,b_1]\).
In the same way \((\sigma'_n(t))_{n\geq 0}\) is a Cauchy sequence in the complete ball \(\overline{B}(\sigma'(t),\varepsilon)\) for all \(t\in [a_1,b]\).
The local geodesics \(\sigma_n\) and \(\sigma'_n\) therefore converge uniformly to local geodesics \(\sigma_\infty:[a,b_1]\to X\) and \(\sigma'_\infty:[a_1,b]\).
These local geodesics restricted to the interval \([a_1,b_1]\) are two local geodesics between \(p_\infty,q_\infty\).
As \(b_1-a_1<A/2<A\), any local geodesic between these two points must be unique as shown before, so \(\sigma_\infty,\sigma_\infty'\) must coincide on \([a_1,b_1]\).
Now these make up a local geodesic from \(\overline{p}\) to \(\overline{q}\) given by
\[\overline{\sigma}=\begin{cases}
\sigma_\infty(t)   &\text{if }t\in [a,b_1]\\
\sigma'_\infty(t)   &\text{if }t\in[a_1,b]\end{cases}.\]
This finishes the proof.
\end{proof}

\end{document}
